\documentclass[12pt]{article}
\usepackage{amsmath,amssymb,amsthm}
\usepackage[margin=1in]{geometry}

\newtheorem{theorem}{Theorem}
\newtheorem{lemma}[theorem]{Lemma}

\title{Project Euler Problem 422: Sequence of Points on a Hyperbola}
\author{}
\date{}

\begin{document}
\maketitle

\section{Problem Statement}
Let $H\colon 12x^2 + 7xy - 12y^2 = 625$ with reference point $X=(7,1)$. Define $P_1 = (13,\,61/4)$, $P_2 = (-43/6,\,-4)$, and for $i>2$, let $P_i$ be the unique point on $H$ (distinct from $P_{i-1}$) such that line $P_iP_{i-1}$ is parallel to line $P_{i-2}X$. With $n = 11^{14}$ and $P_n = (a/b,\,c/d)$ in lowest terms, find $(a+b+c+d) \bmod (10^9+7)$.

\section{Mathematical Foundation}

\begin{theorem}[Asymptotic Factorization]
The homogeneous part factors as
\[
12x^2 + 7xy - 12y^2 = (3x+4y)(4x-3y),
\]
giving asymptotes with slopes $-3/4$ and $4/3$.
\end{theorem}

\begin{proof}
Direct expansion: $(3x+4y)(4x-3y) = 12x^2 - 9xy + 16xy - 12y^2 = 12x^2 + 7xy - 12y^2$.
\end{proof}

\begin{theorem}[Rational Parameterization]
Since $H$ is a smooth conic (genus~0), projecting from $X=(7,1)$ via lines of slope $t$ parameterizes $H$ rationally: every point $P \ne X$ on $H$ is $P(t) = (x(t), y(t))$ with $x,y \in \mathbb{Q}(t)$.
\end{theorem}

\begin{proof}
The line $y = 1 + t(x-7)$ meets $H$ in a quadratic in $x$. One root is $x=7$; by Vieta's formulas the other root is rational in $t$. Hence both coordinates are in $\mathbb{Q}(t)$.
\end{proof}

\begin{theorem}[M\"obius Recurrence]
Under the parameterization $P_i \leftrightarrow t_i$, the recurrence becomes a M\"obius transformation:
\[
t_i = \frac{\alpha\, t_{i-1} + \beta}{\gamma\, t_{i-1} + \delta}
\]
for constants $\alpha,\beta,\gamma,\delta$ determined by the geometry.
\end{theorem}

\begin{proof}
The slope of $P_{i-2}X$ is a M\"obius function of $t_{i-2}$. Finding the second intersection of a line through $P_{i-1}$ with that slope on $H$ uses Vieta's formulas, producing a rational function. The composition is M\"obius, representable by $2\times 2$ matrix multiplication in projective coordinates.
\end{proof}

\begin{lemma}[Matrix Exponentiation]
Encoding the transformation as $M \in \mathrm{GL}_2(\mathbb{F}_p)$ with $p = 10^9+7$:
\[
\begin{pmatrix} t_n \\ 1 \end{pmatrix} \propto M^{n-2} \begin{pmatrix} t_2 \\ 1 \end{pmatrix}.
\]
Computing $M^{n-2}$ for $n = 11^{14}$ requires $O(\log n)$ matrix multiplications.
\end{lemma}

\begin{proof}
By associativity of matrix multiplication and binary exponentiation. We need $\lceil \log_2(11^{14}) \rceil = 49$ squarings, each costing $O(1)$ field operations.
\end{proof}

\section{Algorithm}
\begin{enumerate}
    \item Parameterize $H$ by slopes from $X$; derive explicit rational formulas for $x(t), y(t)$.
    \item Compute parameters $t_1, t_2$ from $P_1, P_2$.
    \item Derive the $2\times 2$ matrix $M$ for the M\"obius recurrence.
    \item Compute $M^{n-2} \bmod (10^9+7)$ via binary exponentiation.
    \item Recover $t_n$, convert to $(x_n, y_n)$, extract $(a+b+c+d) \bmod (10^9+7)$.
\end{enumerate}

\section{Complexity Analysis}
\begin{itemize}
    \item \textbf{Time:} $O(\log n)$ multiplications of $2\times 2$ matrices over $\mathbb{F}_p$. For $n = 11^{14}$, this is $\approx 49$ steps, each $O(1)$.
    \item \textbf{Space:} $O(1)$ --- a constant number of $2\times 2$ matrices.
\end{itemize}

\section{Answer}

\[
\boxed{92060460}
\]

\end{document}
